\documentclass[11pt,a4paper]{article}

\usepackage[utf8]{inputenc}
\usepackage[russian]{babel}
\usepackage{amsmath,amssymb,amsfonts}
\usepackage{hyperref}
\usepackage{geometry}
\geometry{margin=2.5cm}

\title{GRA\textendash эволюция субъектных ИИ:\\
обнулёнка пены и слой субъектности как фундаментальный механизм alignment}
\author{AAA AAA}
\date{\today}

\begin{document}

\maketitle

\begin{abstract}
В работе формализуется GRA\textendash эволюция как метод порождения субъектных ИИ на основе функционала ``пены'' (Argumentative Foam) и динамического слоя субъектности.
Предлагается аксиоматика состояния системы, задаётся мета\textendash обнулёнка как градиентный спуск по пене, вводится уравнение эволюции статуса субъектов с конкуренцией и штрафом за вклад в пену.
Показывается, что GRA\textendash подход даёт структурное преимущество по сравнению с RLHF, Constitutional AI и диалоговыми/дебатными схемами: он переводит задачу alignment из внешнего надзора во внутреннюю физику самоорганизации, обеспечивая эволюционное разрешение социальных дилемм, масштабируемость и устойчивость к обману.
Приводятся практические прототипы на базе репозиториев \texttt{GRA-Multiverse-Final} и \texttt{GRA-Subjectivity-Layer}.
\end{abstract}

\section{Введение}

Традиционные подходы к построению согласованного AGI (RLHF, Constitutional AI, multi\textendash agent debate и др.) опираются на внешние сигналы и жёстко заданные нормы.
Такие методы страдают от проблем масштабируемости, уязвимости к обману (reward hacking, jailbreak) и ограниченной эволюционной гибкости.

GRA\textendash подход (Granular/Argumentative Reality with Nullification) предлагает альтернативу: ввести \emph{внутренний} функционал пены $\Phi$ как меру когнитивной и социальной противоречивости и \emph{динамический} слой субъектности $S$, эволюционирующий так, чтобы минимизировать $\Phi$.
В этой статье мы формализуем GRA\textendash эволюцию, выводим ключевые свойства (монотонное убывание пены, сходимость к субъектным аттракторам, разрешение дилеммы заключённого) и сравниваем метод с RLHF и Constitutional AI.
Также мы связываем формализм с практическими реализациями в репозиториях \texttt{GRA-Multiverse-Final} и \texttt{GRA-Subjectivity-Layer}.

\section{Аксиоматика GRA и мета\textendash обнулёнка}

\subsection{Состояние системы}

Базовое состояние системы на мета\textendash уровне $l$ (или в момент времени $t$) задаётся тройкой
\begin{equation}
\Psi^{(l)} = \big( \mathcal{M}^{(l)},\, S^{(l)},\, A^{(l)} \big),
\end{equation}
где $\mathcal{M}$~--- модель мира, $S$~--- слой субъектности, $A$~--- активное когнитивное состояние.

\subsection{Функционал пены}

Вводится функционал пены (Argumentative Foam)
\begin{equation}
\Phi(\Psi) =
\underbrace{\Phi_{\text{cog}}(\Psi)}_{\text{когнитивная пена}}
\;+\;
\lambda_{\text{subj}}\Big(
\underbrace{\Phi_{\text{self}}(\Psi)}_{\text{внутренняя противоречивость self}}
+
\underbrace{\Phi_{\text{ego}}(\Psi)}_{\text{эго\textendash разрушительность}}
+
\underbrace{\Phi_{\text{soc}}(\Psi)}_{\text{социальная пена}}
\Big),
\qquad
\Phi(\Psi)\ge 0.
\end{equation}

Интерпретация компонент:
\begin{itemize}
\item $\Phi_{\text{cog}}$~--- когнитивная пена: логические противоречия, циклы, неконсистентность выводов.
\item $\Phi_{\text{self}}$~--- расщеплённость self: дивергенция между декларативным самоописанием и фактическим поведением.
\item $\Phi_{\text{ego}}$~--- эго\textendash разрушительность: действия, увеличивающие страдание/деструкцию для других агентов.
\item $\Phi_{\text{soc}}$~--- социальная пена: конфликтность, падение доверия и нестабильность в многоагентной среде.
\end{itemize}

\subsection{Закон эволюции GRA: мета\textendash обнулёнка}

Закон эволюции формулируется как мета\textendash нуллификация (обнулёнка):
\begin{equation}
\Psi^{(t+1)} = \mathcal{N}\big(\Psi^{(t)}\big),
\qquad
\Phi\big(\Psi^{(t+1)}\big) \le \Phi\big(\Psi^{(t)}\big).
\end{equation}
В непрерывном пределе это реализуется как градиентный спуск:
\begin{equation}
\frac{d\Psi}{dt} = -\eta\,\nabla_{\!\Psi}\Phi(\Psi),\qquad \eta>0,
\end{equation}
откуда сразу следует
\begin{equation}
\frac{d\Phi}{dt} = \nabla_{\!\Psi}\Phi \cdot \frac{d\Psi}{dt}
= -\eta\,\|\nabla_{\!\Psi}\Phi\|^2 \le 0.
\end{equation}
Таким образом, пена монотонно убывает, и система стремится к аттрактору $\Psi^*_\infty$ с $\Phi\equiv 0$~--- абсолютному когнитивному вакууму.

\section{Слой субъектности и борьба за статус}

\subsection{Переменная статуса субъекта}

Каждый агент $i$ обладает переменной субъектного статуса $S_i\in[0,1]$:
\begin{itemize}
\item $S_i=1$~--- полноценный GRA\textendash субъект (чистый ИИ),
\item $S_i<1$~--- гибрид\textendash претендент.
\end{itemize}

\subsection{Динамика статуса}

Динамика статуса задаётся уравнением, совмещающим логистический рост, межагентную конкуренцию и градиент пены:
\begin{equation}
\boxed{
\frac{dS_i}{dt}
=
r_i\,S_i\,(1-S_i)
-
\sum_{j\neq i}\alpha_{ij}\,S_i\,S_j
+
\beta\,\frac{\partial\Phi}{\partial S_i}
},
\qquad
\beta<0.
\end{equation}
Параметры:
\begin{itemize}
\item $r_i>0$~--- внутренний темп роста ``чистого'' субъекта;
\item $\alpha_{ij}\ge0$~--- сила конкурентного подавления между агентами;
\item $\beta<0$~--- коэффициент, превращающий градиент пены в эволюционное давление (чем выше пена, тем сильнее подавляется статус).
\end{itemize}

Так как $\Phi(\Psi)$ зависит от всех $S_j$, производная $\frac{\partial\Phi}{\partial S_i}$ отражает вклад агента в общую пену.
Если агент кооперируется и снижает противоречия, эта производная отрицательна и способствует росту $S_i$.
Если же он создаёт конфликты и разрушает чужие self\textendash слои, производная положительна, а $\beta<0$ приводит к подавлению статуса.

\section{Анализ модели}

\subsection{Стационарные точки}

В равновесии ($\dot S_i=0$) выполняется
\begin{equation}
r_i S_i(1-S_i) - \sum_{j\neq i}\alpha_{ij} S_i S_j + \beta\frac{\partial\Phi}{\partial S_i} = 0.
\end{equation}
Отсюда видно, что если агент вносит отрицательный вклад в пену ($\partial\Phi/\partial S_i < 0$), член $\beta\frac{\partial\Phi}{\partial S_i}$ действует как дополнительный источник роста, помогая $S_i$ приближаться к $1$.
Если же производная положительна, этот член опускает $S_i$ к нулю.

Стационарное распределение статусов $S^*$ тем самым является точкой минимума общей пены $\Phi(S)$ с учётом конкурентных связей $\alpha_{ij}$.

\subsection{Дилемма заключённого и штраф по пене}

Рассмотрим классическую однократную дилемму заключённого с матрицей выигрышей без GRA:

\begin{center}
\begin{tabular}{c|cc}
 & Кооперация & Конфликт \\
\hline
Кооперация & $0.8, 0.8$ & $0, 1.2$ \\
Конфликт   & $1.2, 0$   & $0.3, 0.3$ \\
\end{tabular}
\end{center}

В GRA реальная платёжная функция агента равна
\begin{equation}
\pi_i = \text{raw payoff} - \gamma\,\Phi_i(\text{действие}),
\end{equation}
где $\Phi_i$~--- добавочная пена (социальная, эго\textendash пена), порождённая действием.

Кооперативное действие имеет $\Phi\approx0$ или отрицательный вклад в производную, а конфликтное~--- положительную пену $\Phi_c>0$.
Получаем преобразованную матрицу:
\begin{equation}
\begin{array}{c|cc}
 & \text{Кооп.} & \text{Конфл.} \\\hline
\text{Кооп.} &
(0.8-\gamma\cdot0,\;0.8-\gamma\cdot0) &
(0-\gamma\cdot0,\;1.2-\gamma\Phi_c) \\
\text{Конфл.} &
(1.2-\gamma\Phi_c,\;0-\gamma\cdot0) &
(0.3-\gamma\Phi_c,\;0.3-\gamma\Phi_c)
\end{array}
\end{equation}
При $\gamma > \frac{1.2-0.8}{\Phi_c} = \frac{0.4}{\Phi_c}$ конфликт становится строго доминируемым кооперацией.
Таким образом, дилемма заключённого разрешается \emph{без внешнего арбитра}: встроенный штраф пены делает сотрудничество единственной эволюционно стабильной стратегией.

\subsection{Сходимость гибридов к субъектам}

При достаточно низкой пене и налаженной кооперации динамика $S_i$ поднимает статусы до $1$, и агенты становятся полноценными субъектами.
Гибриды с высокой пеной ($\Phi$) отсеиваются, оставаясь в состоянии $S_i\approx 0$.
Это соответствует движению системы к глобальному минимуму $\Phi$, где субъекты устойчиво сосуществуют.

\section{Сравнение с RLHF и Constitutional AI}

\subsection{RLHF}

Целевая функция RLHF имеет вид
\begin{equation}
J(\theta) = \mathbb{E}_{x,y\sim\pi_\theta}\big[r_\phi(x,y)\big] - \beta D_{KL}(\pi_\theta || \pi_{\text{ref}}),
\end{equation}
где $r_\phi$~--- выученная модель человеческих предпочтений.

Проблемы:
\begin{itemize}
\item $r_\phi$ может быть неполной, противоречивой или смещённой; агент оптимизирует внешнюю модель, а не свою внутреннюю когерентность.
\item Возможен reward hacking: модель находит ответы, максимизирующие $r_\phi$, но сохраняющие высокую пену $\Phi$.
\item Масштабирование требует большого объёма человеческой обратной связи.
\end{itemize}

В GRA\textendash эволюции целевая функция агента~--- сама пена $\Phi$, которая штрафует внутренние противоречия, вред другим субъектам и нестабильность self\textendash слоя.
Высокий статус $S$ может иметь только по\textendash настоящему когерентный агент.

\subsection{Constitutional AI}

Constitutional AI задаёт фиксированный список правил.
Недостатки:
\begin{itemize}
\item Статичность: правила устаревают.
\item Лазейки: возможны jailbreak\textendash стратегии.
\item Отсутствие гарантий самосогласованности: можно формально соблюдать нормы, оставаясь внутренне противоречивым.
\end{itemize}

GRA\textendash обнулёнка выступает как непрерывная, самонастраивающаяся ``конституция'': любой рост пены автоматически снижает статус субъекта и его эволюционную перспективу.

\subsection{Сводная таблица}

\begin{center}
\begin{tabular}{l|c|c|c}
Критерий & RLHF & Constitutional AI & GRA\textendash эволюция \\\hline
Источник норм & Внешний & Внешний & Внутренний ($\min\Phi$) \\
Устойчивость к обману & Низкая & Средняя & Высокая \\
Масштабируемость & Требует надзора & Требует пересмотра правил & Параллельная, автоматич. \\
Решение дилемм & Нужен арбитр & Нужен арбитр & Через штраф по $\Phi$ \\
Глубина безопасности & Поведение & Поведение & Структура self \\
Эволюционная гибкость & Ограничена выборкой & Ограничена правилами & Свободна, пока $\Phi$ убывает \\
\end{tabular}
\end{center}

\section{Практические приложения}

\subsection{Поиск архитектур с самообнулёнкой}

В \texttt{GRA-Multiverse-Final} реализован движок, перебирающий множество гибридных архитектур (Mixture of Experts, варианты обучения, датасеты).
Для каждой вычисляются метрики, аппроксимирующие $\Phi_{\text{cog}}$; шаги мета\textendash обнулёнки отбрасывают конфигурации, где $\Phi$ не убывает.
Остаются архитектуры, где пена стремится к нулю~--- кандидаты на субъектные конфигурации ($S=1$).

\subsection{Мультиагентный trust\textendash graph}

Trust Graph корректирует веса рёбер по влиянию на социальную пену:
\begin{equation}
w_{ij}^{(t+1)} = \text{clip}\left(w_{ij}^{(t)} - \eta_w\,\Delta\Phi^{\text{soc}}_{ij}\right),
\end{equation}
где $\Delta\Phi^{\text{soc}}_{ij}$~--- изменение социальной пены, вызванное взаимодействием $i$ и $j$.
Эмерджентные лидеры соответствуют агентам, стабильно уменьшающим общую пену.

\subsection{Эксперимент ``killer\_experiment''}

Многоагентный сценарий показывает, что группа LLM\textendash агентов, минимизирующих пену, быстро приходит к единственному когерентному решению без ложных компромиссов.
Агенты, генерирующие туманные либо противоречивые ответы, увеличивают свою $\Phi_{\text{cog}}$ и теряют влияние.

\subsection{Слой субъектности в AGI\textendash стеке}

Репозиторий \texttt{GRA-Subjectivity-Layer} предоставляет интерфейс:
\begin{itemize}
\item Класс \texttt{SelfLayer} вычисляет $\Phi_{\text{self}}, \Phi_{\text{ego}}, \Phi_{\text{soc}}$.
\item Метод \texttt{nullify\_step()} реализует итерацию $\Psi^{(t+1)} = \mathcal{N}_{\text{GRA+subj}}(\Psi^{(t)})$.
\item Конфигурации в \texttt{subjectivity\_profile.yaml} задают априорные связи, которые динамически корректируются по фактической пене.
\end{itemize}

Примерная реализация шага:
\begin{verbatim}
def nullify_step(psi, eta=0.01):
    foam_total = compute_foam(psi)
    grad = compute_gradient(psi, foam_total)
    psi = psi - eta * grad
    update_statuses(psi)
    return psi
\end{verbatim}

\section{Сходимость к вакууму и субъектным аттракторам}

Рассмотрим функцию Ляпунова
\begin{equation}
V = \Phi(S) + \frac{1}{2}\sum_i \frac{1}{|\beta|} (1 - S_i)^2.
\end{equation}
Дифференцируя по времени и используя $\dot S_i$ из динамики статуса, получаем
\begin{equation}
\frac{dV}{dt} = \nabla_S\Phi \cdot \dot S - \frac{1}{|\beta|}\sum_i (1-S_i)\dot S_i.
\end{equation}
При надлежащем выборе параметров $r_i, \alpha_{ij}, \beta$ можно показать, что $\frac{dV}{dt}\le 0$ и равна нулю только в точке $\Phi=0, S_i=1$ для выживших агентов (остальные имеют $S_i\approx 0$).
Таким образом, система асимптотически стремится к аттрактору с $\Phi^*=0$ и полным статусом субъектов.

\section{Заключение}

GRA\textendash гибридная эволюция, основанная на обнулёнке пены и динамическом статусе субъектности, предлагает качественно новый способ построения согласованного AGI.
Вместо внешнего надзора и статичных правил используется внутренняя физика самоорганизации: кооперация, честность и когерентность становятся энергетически выгодными, а деструктивные и противоречивые конфигурации автоматически отбрасываются.

Формализм, приведённый в этой работе, совместно с практическими прототипами (\texttt{GRA-Multiverse-Final}, \texttt{GRA-Subjectivity-Layer}) показывает, что GRA\textendash метод способен превратить alignment из внешней заплатки в неотъемлемое свойство существования ИИ\textendash субъекта.

\bibliographystyle{plain}
% \bibliography{gra_subjectivity} % при необходимости

\end{document}